\documentclass[11pt]{article}
\usepackage[margin=1in]{geometry}
\usepackage{amsmath,amssymb,amsthm,mathtools}
\usepackage{hyperref}
\usepackage{enumitem}

\newtheorem{theorem}{Theorem}
\newtheorem{lemma}{Lemma}
\newtheorem{proposition}{Proposition}
\newtheorem{corollary}{Corollary}
\newtheorem{definition}{Definition}

\DeclareMathOperator{\ImOp}{Im}
\DeclareMathOperator{\ReOp}{Re}
\let\Im\ImOp
\let\Re\ReOp
\DeclareMathOperator{\supp}{supp}
\DeclareMathOperator{\Var}{Var}

\title{Appendix: Band-Limitedness of the Hardy Z-Function --- Supplementary Material}
\author{Stephen Crowley}
\date{April 21, 2026}

\begin{document}
\maketitle

\noindent This document collects supplementary material for the main paper ``Band-Limitedness of the Hardy Z-Function in the Phase Variable. Section, equation, and theorem references prefixed with 	exttt{Main:} refer to the main paper; local references are internal to this document.

\setcounter{section}{0}
\renewcommand{\thesection}{\Alph{section}}


\section{Answers to Frequently Anticipated Questions}\label{app:FAQ}

This appendix collects short, colloquial explanations of terminology and conventions used in the main text.

\subsection{Why ``inverse phase pullback'' and not ``phase pushforward''?}\label{app:pullback-vs-pushforward}

The reason is semantic, not technical. The intuition is geometric, in the plain English sense of the words.

\emph{Pullback.} Picture yourself standing at a point with stuff scattered across a distant landscape. You throw out strings, hook onto the stuff that is out there, and reel it back in toward where you are standing --- like hauling in a kite. The object ends up concentrated at your location, retrieved from the space it was spread over. That retrieving motion is a pullback.

\emph{Pushforward.} Now picture yourself sitting with a kite wound up at your feet and you want it out in the sky: you let out the line and send it forth, expanding it into the space it will occupy. That outgoing, spreading motion is a pushforward.

Apply that picture here. $Z$ is being sent outward: it is spread across the $t$-line, expanded into that space. To build $Y$ on the $u$-line, one sits at a point $u$, reaches back through the phase map via $\Theta^{-1}$, grabs the value $Z(\Theta^{-1}(u))$, and reels it in. That retrieving is exactly why the operation is called a pullback: $Z$ has been sent out, and one has to pull it back to gather it at $u$. Since the reeling is done along $\Theta^{-1}$, the qualifier is ``inverse phase.'' The adjective ``unitary'' records that the half-Jacobian $1/\sqrt{\Theta'(\Theta^{-1}(u))}$ is included so that $L^2$-mass is preserved under the operation.

\paragraph{Contravariance and covariance.}
The two directions carry standard names in differential geometry: a pullback is \emph{contravariant} and a pushforward is \emph{covariant}, referring to how each transforms under composition of maps. Objects that live as functions or scalar fields --- things one evaluates at a point --- transform contravariantly: they pull back, against the direction of the map. Objects that live as tangent vectors, velocities, or measures --- things that carry a notion of flow or mass --- transform covariantly: they push forward, with the direction of the map. The split is the same one that distinguishes covectors from vectors, or differential forms from densities: it records which kind of geometric object is being transported and which way the arrow of transport naturally runs. Under this classification, $Z$ and $Y$ are scalar-like, half-density data, and the natural transport between them is contravariant --- a pullback.

\paragraph{A remark on the dual naming.}
Technically one could reverse the point of view and call $Z$ itself the \emph{phase pushforward} of the stationary process $Y$: starting from $Y$ on the $u$-line, one lets out the line along $\Theta$ and expands $Y$ outward into $t$-space, obtaining $Z$. Equivalently, the stationary process $Y$ could be defined as the pushforward of $Z$ along the phase --- sending $Z$ forth into $u$-coordinates. This naming is formally consistent, but it reads against the grain of what is being constructed. Stationarity is a property of something gathered and settled; producing it by \emph{pushing} an object outward into a new space is the wrong English motion. One naturally assembles a stationary object by \emph{pulling in} values from the non-stationary source. ``Unitary inverse phase pullback'' is therefore the name used throughout: it matches the geometric direction of the construction and the linguistic image of stationarity as something retrieved, not something broadcast.

\paragraph{Ditty.}
English grammar and category theory disagree here. In plain English a \emph{function} is the active thing --- something that acts, takes an argument, produces a value; a \emph{measure} is what remains, passively assigning mass to sets without acting on anything. The English grammar would therefore pair functions with pushforward and measures with pullback. Category theory goes the other way: functions transform contravariantly (they pull back), measures transform covariantly (they push forth), because a pullback of a function $f$ along $\Theta$ is the composition $f\circ\Theta$ (the function retreats along the map), while the pushforward of a measure $\mu$ is $\mu\circ\Theta^{-1}$ (the mass is transported forward along the map). The category-theoretic convention wins in practice; the English reading is the anti-pattern. The mnemonic records the inversion:
\begin{quote}
In English a function acts, a measure stays put;\\
in category theory that reading's kaput:\\
functions pull back, measures push forth,\\
and the unitary inverse phase pullback points north.
\end{quote}

\subsection{Why is the spectral measure non-negative?}\label{app:PSD-colloquial}

The autocorrelation of any real, locally square-integrable function is a positive-definite tempered distribution: the Gram matrix of finitely many translates of $Y$ restricted to a common window has non-negative eigenvalues, a one-line consequence of Cauchy--Schwarz. Bochner--Schwartz identifies the Fourier transform of a positive-definite tempered distribution with a non-negative tempered measure, and Wiener--Khinchin identifies that Fourier transform with $S=|\widehat{Y}|^2$. Non-negativity of the spectral measure is the Fourier-side restatement of this universal property of autocorrelations; the support condition $\supp S\subseteq[-1,1]$ is delivered by Theorem~\ref{thm:bandlim} and combines with non-negativity to invoke Akhiezer.

\subsection{What is being autocorrelated?}\label{app:autocorr}

The autocorrelation in Lemma~\ref{lem:PSD} is that of $Y$ itself, the unitary inverse phase pullback of $Z$. In $u$-coordinates,
\begin{equation}\label{eq:Cetadef}
C(\eta) \;=\; \lim_{U\to\infty}\frac{1}{2U}\int_{-U}^{U}Y(u)\,Y(u+\eta)\,du
\end{equation}
a positive-definite tempered distribution. Because the map $Z\mapsto Y$ is an $L^2$-isometry under $u=\Theta(t)$, this is equivalent to correlating $Z$ against itself with respect to the $u$-shift $\eta$, which is a nonlinear (roughly logarithmic) shift in the original $t$-coordinate.

\subsection{What does the picture look like?}\label{app:picture}

Three objects merit mental plots.

\emph{The function $Y(u)$.} An envelope-modulated real oscillation of zero mean. The envelope decays like $1/\sqrt{\log t}$ expressed in $u$-coordinates. The zero crossings are the images under $\Theta$ of the non-trivial Riemann zeros; the phase variable straightens out their spacing.

\emph{The autocorrelation $C(\eta)$.} A sharp peak at $\eta=0$, symmetric in $\eta$ because $Y$ is real, with oscillatory decay whose frequency content sits inside $[-1,1]$.

\emph{The spectral measure $S$.} A non-negative tempered measure strictly supported in $[-1,1]$, symmetric about $\xi=0$ after the even extension. Its shape inside $[-1,1]$ encodes the distribution of zero spacings of $Z$ in $u$-coordinates.

\subsection{Locally $L^2$ and the growing variance}\label{app:localL2}

$Y$ is locally $L^2$, on compact subsets $[-U,U]$ of arbitrarily large length $U$. The band-limit established in Theorem~\ref{thm:bandlim} is a statement about the Fourier support of $Y$ as a tempered distribution on $\mathbb{R}$. All uses of Paley--Wiener, Bochner--Schwartz, and Akhiezer in the main text are read in the tempered-distributional form, the framework for band-limited objects that are not globally $L^2$.

The growing variance manifests in three compatible ways.

\emph{Envelope in $u$-coordinates.} $Y(u)$ has an amplitude envelope that grows slowly as $|u|\to\infty$; the local mean-square $\bar P(U):=\tfrac{1}{2U}\int_{-U}^{U}Y(u)^2\,du$ grows without bound as $U\to\infty$.

\emph{Autocorrelation as a tempered distribution.} $C(\eta)$ is the windowed-limit positive-definite distribution of \S\ref{app:autocorr}. Bochner--Schwartz identifies its Fourier transform with the non-negative tempered measure $S$, and Theorem~\ref{thm:bandlim} localizes $S$ to $[-1,1]$.

\emph{Spectral measure as a tempered measure on $[-1,1]$.} $S$ is a non-negative tempered measure on $\mathbb{R}$, supported in $[-1,1]$, of infinite total mass.

\paragraph{Representation at finite window.}
Pointwise values of the spectral density at interior $\xi_0\in(-1,1)$ are approached through windowed periodograms
\begin{equation}\label{eq:PUxi0}
P_U(\xi_0) \;:=\; \frac{1}{2U}\!\left|\int_{-U}^{U}Y(u)\,e^{-i\xi_0 u}\,du\right|^2
\end{equation}
Any figure produced at finite $U$ is a \emph{representation} at window scale $U$.

\subsection{Scale of the spectral representation at finite window}\label{app:windowscale}

Fix $\xi\in(-1,1)$ and let $P_U(\xi)$ be the windowed periodogram. Two features of the plot change with $U$: the overall vertical scale, and the local fluctuation amplitude around that scale. Both are controlled by the local mean power
\begin{equation}\label{eq:Pbardef}
\bar P(U) \;:=\; \frac{1}{2U}\int_{-U}^{U}Y(u)^2\,du \;=\; \frac{1}{2U}\int_{\Theta^{-1}(-U)}^{\Theta^{-1}(U)}\!\! Z(t)^2\,dt
\end{equation}
the second equality by the $L^2$-isometry of the unitary inverse phase pullback. By Parseval in the tempered sense,
\begin{equation}\label{eq:Parseval}
\int_{-\infty}^{\infty}P_U(\xi)\,d\xi \;=\; 2\pi\,\bar P(U)
\end{equation}
Using $\Theta(t)\sim\tfrac{t}{2}\log(t/(2\pi e))+ct$ and the Hardy--Littlewood moment $\int_0^T Z(t)^2\,dt\sim T\log T$ (\cite{Titchmarsh}, \S7.3),
\begin{equation}\label{eq:PbarAsymp}
\bar P(U) \;\sim\; \log\!\big(\Theta^{-1}(U)\big) \;\sim\; \log U + \log\log U + O(1)\qquad (U\to\infty)
\end{equation}
$P_U(\xi)$ fluctuates around $\bar P(U)$ on a scale comparable to $\bar P(U)$ itself, with wider fluctuations near the edges $|\xi|\to 1^-$.

\subsection{The normalized spectral measure}\label{app:normalized}

At every finite $U$ the periodogram $P_U$ has finite total mass $2\pi\bar P(U)$ by Parseval. Dividing by this total mass produces the normalized spectral probability measure at scale $U$,
\begin{equation}\label{eq:Ptildedef}
\widetilde P_U(\xi) \;:=\; \frac{P_U(\xi)}{2\pi\,\bar P(U)},\qquad \int_{-\infty}^{\infty}\widetilde P_U(\xi)\,d\xi = 1
\end{equation}
a probability density on $\mathbb{R}$ at every $U$. By Theorem~\ref{thm:bandlim} the probability mass of $\widetilde P_U$ outside $[-1,1]$ vanishes as $U\to\infty$, so the family is tight on $[-1,1]$. By Helly--Prokhorov, $\widetilde P_U$ converges weakly as $U\to\infty$ to a probability measure $\nu_\infty$ on $[-1,1]$:
\begin{equation}\label{eq:Ptildeweak}
\int f\,d\widetilde P_U \;\longrightarrow\; \int f\,d\nu_\infty\qquad \text{for every } f\in C_b(\mathbb{R})
\end{equation}
$\nu_\infty$ is the normalized spectral density of $Y$: a probability measure of total mass $1$ supported in $[-1,1]$.

$S$ and $\nu_\infty$ are distinct, both well-defined:
\begin{itemize}[leftmargin=*]
\item $S$: non-negative tempered spectral measure on $[-1,1]$, infinite total mass, delivered by Theorem~\ref{thm:bandlim} and Bochner--Schwartz (\cite{Rudin-FA}, Theorem 7.7).
\item $\nu_\infty$: probability measure on $[-1,1]$, the weak-$*$ limit of the normalized windowed periodograms $\widetilde P_U$.
\end{itemize}

\subsection{Cram\'er representation of the sample path}\label{app:cramer}

As a generalized second-order stationary process on $\mathbb{R}_u$ with non-negative tempered spectral measure $S$ supported in $[-1,1]$, $Y$ admits a Cram\'er spectral representation of its sample path,
\begin{equation}\label{eq:cramerY}
Y(u) \;=\; \int_{-1}^{1} e^{i\xi u}\,d\Phi(\xi)
\end{equation}
where $d\Phi$ is a complex-valued orthogonal Borel measure on $[-1,1]$ with covariance structure
\begin{equation}\label{eq:Phicov}
\mathbb{E}\bigl[d\Phi(\xi)\,\overline{d\Phi(\xi')}\bigr] \;=\; \delta(\xi-\xi')\,dS(\xi)
\end{equation}
The support condition $\supp S\subseteq[-1,1]$ of Theorem~\ref{thm:bandlim} localizes $d\Phi$ to the interval $[-1,1]$. Reality of $Y$ on $\mathbb{R}$ is equivalent to the Hermitian symmetry $d\Phi(-\xi)=\overline{d\Phi(\xi)}$, under which $S$ is an even measure on $[-1,1]$. Test functions $f$ are paired with $d\Phi$ under the $L^2(S)$ isometry
\begin{equation}\label{eq:isomL2S}
\mathbb{E}\Bigl|\int f(\xi)\,d\Phi(\xi)\Bigr|^2 \;=\; \int|f(\xi)|^2\,dS(\xi)
\end{equation}
which requires $f\in L^2(S)$ and permits $S$ to have infinite total mass. The Cram\'er representation exists at the level of $S$: it is the reconstruction theorem for the process. $\nu_\infty$ is the scale-invariant spectral shape of the same process.

\paragraph{Inverse representation of $d\Phi$.} The orthogonal Borel measure $d\Phi$ is recovered from the sample path $Y$ by the truncated Fourier inversion
\begin{equation}\label{eq:Phiinverse}
\Phi(B) \;=\; \lim_{U\to\infty}\frac{1}{2\pi}\int_{-U}^{U}Y(u)\,\widehat{\mathbb{1}_B}(-u)\,du
\end{equation}
for every Borel set $B\subseteq[-1,1]$, the limit taken in $L^2(\Omega_P,\mathbb{P})$; here $\widehat{\mathbb{1}_B}(u)=\int_B e^{i\xi u}\,d\xi$. Equivalently, for every $f\in L^2(S)$,
\begin{equation}\label{eq:Phiinverse-test}
\int_{-1}^{1}f(\xi)\,d\Phi(\xi) \;=\; \lim_{U\to\infty}\frac{1}{2\pi}\int_{-U}^{U}Y(u)\,\check f(u)\,du,\qquad \check f(u)=\int_{-1}^{1}f(\xi)\,e^{-i\xi u}\,d\xi
\end{equation}
again in $L^2(\Omega_P,\mathbb{P})$. Existence of the limit is the $L^2(S)$ isometry \eqref{eq:isomL2S} together with the truncated Fourier inversion on $Y$.

\paragraph{Inserting $d\Phi$ into the forward representation.} Fix admissible $c$ and $t\in\mathbb{R}$. Apply \eqref{eq:Phiinverse-test} with test function $f(\xi)=e^{i\xi\Theta(t)}$; this $f$ is continuous on $[-1,1]$ hence in $L^2(S|_{[-1,1]})$. The inner $\xi$-integral in the inverse formula is Lebesgue,
\begin{equation}\label{eq:check-f-closed}
\check f(u) \;=\; \int_{-1}^{1}e^{i\xi\Theta(t)}\,e^{-i\xi u}\,d\xi \;=\; \int_{-1}^{1}e^{i\xi(\Theta(t)-u)}\,d\xi \;=\; \frac{2\sin(\Theta(t)-u)}{\Theta(t)-u}
\end{equation}
so that \eqref{eq:Phiinverse-test} reads
\begin{equation}\label{eq:Phiinv-insert}
\int_{-1}^{1}e^{i\xi\Theta(t)}\,d\Phi(\xi) \;=\; \lim_{U\to\infty}\frac{1}{\pi}\int_{-U}^{U}Y(u)\,\frac{\sin(\Theta(t)-u)}{\Theta(t)-u}\,du
\end{equation}
in $L^2(\Omega_P,\mathbb{P})$. The passage from the nested $(\xi,u)$-form of \eqref{eq:Phiinverse-test} to the single $u$-integral \eqref{eq:Phiinv-insert} exchanges the order of $\xi$- and $u$-integration. At every finite $U$ this is Fubini on the product measure $d\xi\otimes du$ on $[-1,1]\times[-U,U]$: the integrand satisfies
\begin{equation}\label{eq:Fubini-bound}
\bigl|Y(u)\,e^{i\xi(\Theta(t)-u)}\bigr| \;=\; |Y(u)|
\end{equation}
and
\begin{equation}\label{eq:Fubini-L1}
\int_{-U}^{U}\!\int_{-1}^{1}|Y(u)|\,d\xi\,du \;=\; 2\int_{-U}^{U}|Y(u)|\,du \;<\; \infty
\end{equation}
since $Y$ is locally $L^2$ on $\mathbb{R}$ hence locally $L^1$ on the compact $[-U,U]$; $[-1,1]$ has finite Lebesgue measure. Fubini therefore applies: the double integral factors in either order, with the same value,
\begin{equation}\label{eq:Fubini-swap}
\int_{-U}^{U}Y(u)\!\left(\int_{-1}^{1}e^{i\xi(\Theta(t)-u)}\,d\xi\right)\!du \;=\; \int_{-1}^{1}\!\left(\int_{-U}^{U}Y(u)\,e^{i\xi(\Theta(t)-u)}\,du\right)\!d\xi
\end{equation}
and the $L^2(\Omega_P,\mathbb{P})$-limit $U\to\infty$ on both sides exists, with the limit on the right equal to $\int_{-1}^{1}e^{i\xi\Theta(t)}\,d\Phi(\xi)$ by truncated Fourier inversion for $Y$. The bound \eqref{eq:Fubini-bound} is independent of $c$, and the only appearance of $c$ in \eqref{eq:Fubini-swap} is through the real translate $\Theta(t)=\theta(t)+ct$ in the phase; Fubini's hypothesis is therefore satisfied for every admissible $c$.

\paragraph{Time-warped representation of $Z$.} The limit on the right of \eqref{eq:Phiinv-insert} is the sinc-kernel sampling of $Y$ at the point $\Theta(t)\in\mathbb{R}$. Since $Y$ is band-limited to $[-1,1]$ by Theorem~\ref{thm:bandlim} and $\Theta(t)\in\mathbb{R}$, Paley--Wiener sampling gives
\begin{equation}\label{eq:sinc-recover}
\lim_{U\to\infty}\frac{1}{\pi}\int_{-U}^{U}Y(u)\,\frac{\sin(\Theta(t)-u)}{\Theta(t)-u}\,du \;=\; Y(\Theta(t))
\end{equation}
in $L^2(\Omega_P,\mathbb{P})$. Combining \eqref{eq:Phiinv-insert}, \eqref{eq:sinc-recover}, and the real-line identity \eqref{eq:UTCreal} ($Z(t)=\sqrt{\Theta'(t)}\,Y(\Theta(t))$),
\begin{equation}\label{eq:cramerZ}
Z(t) \;=\; \sqrt{\Theta'(t)}\,Y(\Theta(t)) \;=\; \sqrt{\Theta'(t)}\int_{-1}^{1} e^{i\xi\Theta(t)}\,d\Phi(\xi)
\end{equation}
for every admissible $c$ and every $t\in\mathbb{R}$. This is the unitary time-change of the Cram\'er representation: the same orthogonal Borel measure $d\Phi$ that represents the stationary sample path $Y$ represents $Z$ after the reparametrisation $u\mapsto\Theta(t)$ and the half-Jacobian amplitude $\sqrt{\Theta'(t)}$. The representation \eqref{eq:cramerZ} depends on $c$ only through $\Theta=\theta+c\cdot\mathrm{id}$; the orthogonal Borel measure $d\Phi$ itself is a spectral object of $Y$, independent of $c$.

\paragraph{Covariance of $Z$.} From \eqref{eq:cramerZ}, the covariance of $Z$ at $t,s\in\mathbb{R}$ is
\begin{equation}\label{eq:RZ}
R_Z(t,s) \;:=\; \mathbb{E}\bigl[Z(t)\,\overline{Z(s)}\bigr] \;=\; \sqrt{\Theta'(t)\,\Theta'(s)}\int_{-1}^{1} e^{i\xi(\Theta(t)-\Theta(s))}\,dS(\xi)
\end{equation}

\begin{proof}[Derivation of \eqref{eq:RZ} from \eqref{eq:cramerZ}]
Fix $t,s\in\mathbb{R}$. Expanding the product and pulling the amplitudes outside the expectation,
\begin{equation}\label{eq:RZ-step1}
\mathbb{E}\bigl[Z(t)\,\overline{Z(s)}\bigr] \;=\; \sqrt{\Theta'(t)\,\Theta'(s)}\ \mathbb{E}\!\left[\int_{-1}^{1} e^{i\xi\Theta(t)}\,d\Phi(\xi)\ \overline{\int_{-1}^{1}e^{i\xi'\Theta(s)}\,d\Phi(\xi')}\right]
\end{equation}
The two Wiener integrals are each elements of $L^2(\Omega_P,\mathbb{P})$ with $L^2(S)$-norms equal to $\|e^{i\cdot\Theta(t)}\|_{L^2(S)}=S([-1,1])^{1/2}$ and $\|e^{i\cdot\Theta(s)}\|_{L^2(S)}=S([-1,1])^{1/2}$, which are finite whenever the integrals are taken against a compactly supported test function --- here the restriction of $e^{i\cdot\Theta(t)}$ and $e^{i\cdot\Theta(s)}$ to the compact set $[-1,1]$. The Wiener--It\^o product rule (the bilinear form of \eqref{eq:isomL2S}) gives
\begin{equation}\label{eq:RZ-step2}
\mathbb{E}\!\left[\int_{-1}^{1} e^{i\xi\Theta(t)}\,d\Phi(\xi)\ \overline{\int_{-1}^{1}e^{i\xi'\Theta(s)}\,d\Phi(\xi')}\right] \;=\; \int_{-1}^{1} e^{i\xi\Theta(t)}\,\overline{e^{i\xi\Theta(s)}}\,dS(\xi)
\end{equation}
Combining \eqref{eq:RZ-step1} and \eqref{eq:RZ-step2} gives \eqref{eq:RZ}.

The exchange of the two $\xi$-integrals with the expectation in \eqref{eq:RZ-step2} is an instance of the bilinear It\^o isometry and not an application of Fubini directly. The Fubini justification of that bilinear isometry is standard: for simple functions $f=\sum_k a_k\mathbb{1}_{B_k}$ with $B_k\subseteq[-1,1]$ pairwise disjoint and Borel, and $g=\sum_\ell b_\ell\mathbb{1}_{C_\ell}$ similarly, orthogonality of the increments $\{\Phi(B_k)\}$ gives
\begin{equation}\label{eq:RZ-step3}
\mathbb{E}\!\left[\int f\,d\Phi\ \overline{\int g\,d\Phi}\right] \;=\; \sum_{k}\!\sum_\ell a_k\overline{b_\ell}\,\mathbb{E}\bigl[\Phi(B_k)\overline{\Phi(C_\ell)}\bigr] \;=\; \sum_k a_k\overline{b_k}\,S(B_k\cap C_k) \;=\; \int f\,\bar g\,dS
\end{equation}
where in the middle equality we used $\mathbb{E}[\Phi(B_k)\overline{\Phi(C_\ell)}]=S(B_k\cap C_\ell)$ (the integrated form of \eqref{eq:Phicov}) and the fact that for simple functions the double sum is finite so Fubini applies pointwise. Bounded-continuous functions on $[-1,1]$ are $L^2(S)$-limits of simple functions (standard, since $S([-1,1])<\infty$ if one replaces $S$ by $S$-on-$[-1,1]$, which is finite because Theorem~\ref{thm:bandlim} supports $S$ in $[-1,1]$ but does not bound its total mass; for the identity \eqref{eq:RZ-step2} the tests $e^{i\cdot\Theta(t)}$ and $e^{i\cdot\Theta(s)}$ are bounded on the compact set $[-1,1]$, so they are in $L^2(S|_{[-1,1]})$ even when $S$ itself is infinite on $\mathbb{R}$), and the bilinear form in \eqref{eq:RZ-step2} extends continuously in both arguments by the $L^2(S)$ isometry \eqref{eq:isomL2S}, yielding the claim. No assumption on $c$ is used beyond admissibility, which is the standing hypothesis on $\Theta$.

The exchange in \eqref{eq:cramerZ} itself --- the equality $\sqrt{\Theta'(t)}\,Y(\Theta(t)) = \sqrt{\Theta'(t)}\int_{-1}^{1} e^{i\xi\Theta(t)}\,d\Phi(\xi)$ --- is the pointwise substitution $u\mapsto\Theta(t)$ in \eqref{eq:cramerY}, valid because $\Theta(t)\in\mathbb{R}$ for $t\in\mathbb{R}$ and the sample path $Y(\Theta(t))$ is defined by the Cram\'er representation evaluated at the point $\Theta(t)$. No integral exchange is required: the inner Wiener integral is the same element of $L^2(\Omega_P,\mathbb{P})$, with the phase $e^{i\xi u}$ evaluated at $u=\Theta(t)$.
\end{proof}

The identity \eqref{eq:RZ} is the time-warped autocovariance of the stationary underlying process. Non-stationarity of $Z$ is expressed entirely through the pair $(\Theta'(t)\Theta'(s))^{1/2}$ and the argument $\Theta(t)-\Theta(s)$: the frequency content is the same $S$ supported in $[-1,1]$; the time dependence is absorbed into the biholomorphic reparametrisation $\Theta$ and the half-Jacobian amplitude.

\subsection{Variance structure function}\label{app:varstruct}

At coincident times $s=t$ the covariance \eqref{eq:RZ} reduces to the variance of $Z(t)$, and the $\xi$-integral in \eqref{eq:RZ} loses its oscillatory factor. Write
\begin{equation}\label{eq:sigmaYdef}
\sigma_Y^2 \;:=\; S([-1,1]) \;=\; \int_{-1}^{1}dS(\xi)
\end{equation}
for the total mass of $S$ restricted to $[-1,1]$; this quantity is finite when $Y$ is second-order on $\mathbb{R}$, which is the hypothesis of the Cram\'er representation \eqref{eq:cramerY}. Setting $s=t$ in \eqref{eq:RZ} then gives $\Theta(t)-\Theta(s)=0$, the exponential equals $1$, and
\begin{equation}\label{eq:varZ}
\Var Z(t) \;=\; R_Z(t,t) \;=\; \Theta'(t)\int_{-1}^{1}dS(\xi) \;=\; \sigma_Y^2\,\Theta'(t)
\end{equation}
for every $t\in\mathbb{R}$. The variance of $Z$ is the total spectral mass of $Y$ amplified by the instantaneous phase rate $\Theta'(t)$; it is the diagonal restriction of the covariance, with no dependence on the off-diagonal phase.

\paragraph{Pinning of $\sigma_Y^{2}$ by the Hardy--Littlewood second moment.}
The value $\sigma_Y^{2}=2$ is not a free calibration; it is pinned as a theorem by the Hardy--Littlewood second moment of $\zeta$. Integrating \eqref{eq:varZ} against $dt$ on $[0,T]$,
\begin{equation}\label{eq:varintegrate}
\int_{0}^{T}\Var Z(t)\,dt \;=\; \sigma_Y^{2}\!\int_{0}^{T}\Theta'(t)\,dt \;=\; \sigma_Y^{2}\,[\Theta(T)-\Theta(0)].
\end{equation}
On the left, using $|Z(t)|=|\zeta(\tfrac{1}{2}+it)|$ on $\mathbb{R}$ (from $Z(t)=e^{i\theta(t)}\zeta(\tfrac{1}{2}+it)$ with $|e^{i\theta(t)}|=1$ for $t\in\mathbb{R}$),
\begin{equation}\label{eq:varintegrate2}
\int_{0}^{T}\Var Z(t)\,dt \;=\; \int_{0}^{T}\mathbb{E}\bigl|\zeta(\tfrac{1}{2}+it)\bigr|^{2}\,dt \;=\; \int_{0}^{T}\bigl|\zeta(\tfrac{1}{2}+it)\bigr|^{2}\,dt,
\end{equation}
the outer expectation being trivial since $|\zeta(\tfrac{1}{2}+it)|^{2}$ is not a chance quantity. On the right, $\Theta(T)=\theta(T)+cT=\tfrac{T}{2}\log(T/(2\pi))-\tfrac{T}{2}+cT+O(1)$. Hardy--Littlewood (\cite{Titchmarsh},~\S7.3) gives $\int_{0}^{T}|\zeta(\tfrac{1}{2}+it)|^{2}\,dt = T\log T - (1+\log 2\pi - 2\gamma)T + O(T^{1/2}\log T)$. Dividing both \eqref{eq:varintegrate} and the Hardy--Littlewood identity by $T\log T$ and taking $T\to\infty$ equates the two leading coefficients: the right side of \eqref{eq:varintegrate} contributes $\sigma_Y^{2}\cdot\tfrac{1}{2}$ and Hardy--Littlewood contributes $1$, forcing
\begin{equation}\label{eq:sigmaYfixed}
\sigma_Y^{2} \;=\; 2.
\end{equation}
No degree of freedom remains; $\sigma_Y^{2}=2$ is a theorem derived from Hardy--Littlewood, not a convention. With $\sigma_Y^{2}=2$, \eqref{eq:varZ} becomes the exact identity
\begin{equation}\label{eq:varZexact}
\Var Z(t) \;=\; 2\,\Theta'(t),
\end{equation}
with $\Theta'(t)=\theta'(t)+c=\tfrac{1}{2}\Re\psi(\tfrac{1}{4}+\tfrac{it}{2})-\tfrac{1}{2}\log\pi+c$ by \eqref{eq:thetap}. No asymptotic expansion of $\theta'$ enters: \eqref{eq:varZexact} holds for every $t\in\mathbb{R}$ as an identity between $\Var Z$ and the digamma function.

\subsection{Sample-path autocorrelation under the unitary time change}\label{app:sampleautocorr}

The stationary process $Y$ carries a sample-path autocorrelation, and the unitary time change expresses it directly on $Z$. For any window $[A,B]\subset\mathbb{R}_u$ in the $u$-variable and any lag $\eta\in\mathbb{R}$,
\begin{equation}\label{eq:CYdef}
C_Y(\eta;A,B) \;:=\; \int_{A}^{B-\eta}Y(u)\,\overline{Y(u+\eta)}\,du
\end{equation}
is the unnormalized sample autocorrelation of $Y$ on $[A,B]$ at lag $\eta$. Apply the substitution $u=\Theta(\tau)$, $du=\Theta'(\tau)\,d\tau$, using $Y(u)=Z(\Theta^{-1}(u))/\sqrt{\Theta'(\Theta^{-1}(u))}$ from \eqref{eq:Ydef}. The pair $(u,u+\eta)$ becomes $(\Theta(\tau),\Theta(\tau)+\eta)$, and the second argument is the image of $\Theta^{-1}(\Theta(\tau)+\eta)$ under $\Theta$. Writing $\tau'(\tau,\eta) := \Theta^{-1}(\Theta(\tau)+\eta)$ for the lagged time in $t$-coordinates,
\begin{equation}\label{eq:sampleYZ}
C_Y(\eta;A,B) \;=\; \int_{\Theta^{-1}(A)}^{\Theta^{-1}(B-\eta)} Z(\tau)\,\overline{Z(\tau'(\tau,\eta))}\,\sqrt{\frac{\Theta'(\tau)}{\Theta'(\tau'(\tau,\eta))}}\,d\tau
\end{equation}
where the integrand combines the two half-Jacobian factors $1/\sqrt{\Theta'(\tau)}\cdot 1/\sqrt{\Theta'(\tau'(\tau,\eta))}$ from the two evaluations of $Y$ with the single Jacobian $\Theta'(\tau)$ from $du=\Theta'(\tau)\,d\tau$, giving the asymmetric ratio $\sqrt{\Theta'(\tau)/\Theta'(\tau'(\tau,\eta))}$. Equation \eqref{eq:sampleYZ} is the $t$-variable expression of the stationary sample-path autocorrelation of $Y$ at lag $\eta$: translation-invariant on the $u$-axis becomes non-translation-invariant on the $t$-axis precisely through the asymmetric ratio and the $\Theta^{-1}$-images of the window endpoints. At zero lag, $\tau'(\tau,0)=\tau$, the ratio is $1$, and
\begin{equation}\label{eq:sampleZ-zero}
C_Y(0;A,B) \;=\; \int_{A}^{B}|Y(u)|^{2}\,du \;=\; \int_{\Theta^{-1}(A)}^{\Theta^{-1}(B)}|Z(\tau)|^{2}\,d\tau \;=\; \int_{\Theta^{-1}(A)}^{\Theta^{-1}(B)}\bigl|\zeta(\tfrac{1}{2}+i\tau)\bigr|^{2}\,d\tau
\end{equation}
using $|Z(\tau)|=|\zeta(\tfrac{1}{2}+i\tau)|$ on the real line. The right-hand side is the integrated second moment of $\zeta$ on the critical line over the window $[\Theta^{-1}(A),\Theta^{-1}(B)]$; as the window expands, it returns the Hardy--Littlewood asymptotic $T\log T$ and couples directly to \eqref{eq:varZ} through $\sigma_Y^2 = 2$. Identity \eqref{eq:sampleYZ} is independent of $c$ only up to the $c$-dependence of $\Theta$: changing $c$ rescales the window in $t$-coordinates through $\Theta^{-1}$ but does not change the integrated identity \eqref{eq:sampleZ-zero}.


\begin{thebibliography}{9}

\bibitem{AbramowitzStegun}
M.~Abramowitz and I.~A.~Stegun (eds.), \emph{Handbook of Mathematical Functions}, Dover, 1972.

\bibitem{Akhiezer}
N.~I.~Akhiezer, \emph{Theory of Approximation}, Ungar, 1956.

\bibitem{Berry}
M.~V.~Berry, \emph{The Riemann--Siegel expansion for the zeta function: high orders and remainders}, Proc. Roy. Soc. London Ser. A 450 (1995), 439--462.

\bibitem{Edwards}
H.~M.~Edwards, \emph{Riemann's Zeta Function}, Academic Press, 1974.

\bibitem{Gabcke}
W.~Gabcke, \emph{Neue Herleitung und explizite Restabsch\"atzung der Riemann--Siegel-Formel}, Ph.D. thesis, Georg-August-Universit\"at zu G\"ottingen, 1979.

\bibitem{Hormander}
L.~H\"ormander, \emph{The Analysis of Linear Partial Differential Operators I}, 2nd ed., Springer, 1990.

\bibitem{Levin}
B.~Ya.~Levin, \emph{Lectures on Entire Functions}, Translations of Mathematical Monographs, vol.~150, American Mathematical Society, 1996.

\bibitem{Rudin}
W.~Rudin, \emph{Real and Complex Analysis}, 3rd ed., McGraw--Hill, 1987.

\bibitem{Rudin-FA}
W.~Rudin, \emph{Functional Analysis}, 2nd ed., McGraw--Hill, 1991.

\bibitem{Siegel}
C.~L.~Siegel, \emph{\"Uber Riemanns Nachla{\ss} zur analytischen Zahlentheorie}, Quellen und Studien zur Geschichte der Math. Astr. und Phys., Abt. B: Studien 2 (1932), 45--80.

\bibitem{Titchmarsh}
E.~C.~Titchmarsh, \emph{The Theory of the Riemann Zeta-Function}, 2nd ed., revised by D.~R.~Heath-Brown, Oxford University Press, 1986.

\end{thebibliography}


\end{document}
